

% ---
% Configurações de aparência do PDF final

% alterando o aspecto da cor azul
\definecolor{blue}{RGB}{41,5,195}

% informações do PDF
\makeatletter
\hypersetup{
     	%pagebackref=true,
		pdftitle={\@title}, 
		pdfauthor={\@author},
    	pdfsubject={\imprimirpreambulo},
	    pdfcreator={\@author},
		pdfkeywords={palavra-chave}{palavra-chave}{palavra-chave}, 
		colorlinks=true,       		% false: boxed links; true: colored links
    	linkcolor=black,          	% color of internal links
    	citecolor=black,        		% color of links to bibliography
    	filecolor=magenta,      		% color of file links
		urlcolor=black,
		bookmarksdepth=4}
\makeatother
% --- 

% ---
% Posiciona figuras e tabelas no topo da página quando adicionadas sozinhas
% em um página em branco. Ver https://github.com/abntex/abntex2/issues/170
\makeatletter
\setlength{\@fptop}{5pt} % Set distance from top of page to first float
\makeatother
% ---

% ---
% Possibilita criação de Quadros e Lista de quadros.
% Ver https://github.com/abntex/abntex2/issues/176
%
\newcommand{\quadroname}{Quadro}
\newcommand{\listofquadrosname}{Lista de quadros}

\newfloat[chapter]{quadro}{loq}{\quadroname}
\newlistof{listofquadros}{loq}{\listofquadrosname}
\newlistentry{quadro}{loq}{0}

% configurações para atender às regras da ABNT
\setfloatadjustment{quadro}{\centering}
\counterwithout{quadro}{chapter}
\renewcommand{\cftquadroname}{\quadroname\space} 
\renewcommand*{\cftquadroaftersnum}{\hfill--\hfill}

\setfloatlocations{quadro}{hbtp} % Ver https://github.com/abntex/abntex2/issues/176
% ---

% Float de código, no mesmo molde do "quadro" acima:
% legenda em cima, fonte embaixo, numeração contínua e lista própria.
\newcommand{\codigoname}{Código}
\newcommand{\listofcodigosname}{Lista de códigos}
\newfloat[chapter]{codigo}{loc}{\codigoname}
\newlistof{listofcodigos}{loc}{\listofcodigosname}
\newlistentry{codigo}{loc}{0}
\setfloatadjustment{codigo}{\centering}
\counterwithout{codigo}{chapter}
\renewcommand{\cftcodigoname}{\codigoname\space}
\renewcommand*{\cftcodigoaftersnum}{\hfill--\hfill}
\setfloatlocations{codigo}{htbp}

% ---

% --- 
% Espaçamentos entre linhas e parágrafos 
% --- 

% O tamanho do parágrafo é dado por:
\setlength{\parindent}{1.2cm}
% Controle do espaçamento entre um parágrafo e outro:
\setlength{\parskip}{0.3cm}  % tente também \onelineskip
\linespread{1.7}
\renewcommand{\baselinestretch}{1.5}
%\DeclareUnicodeCharacter{0301}{i}

\setlength\bibitemsep{4\itemsep}
